Before clustering, the sampled workspace is restricted to points above $z = 115$~mm for the single-section arm and $z = 200$~mm for the multi-section arm (89.9\% and 99.7\% of the samples), and each cluster center is replaced by the nearest sampled point. The resulting targets are listed in Tables~S8 and S9. The convex hull of the 20 targets covers 9.2\% (single section) and 16.4\% (multi-section) of the volume of the convex hull of the sampled positions, and the median distance from a sampled position to the nearest target is 7.6~mm and 38.5~mm, respectively. For the single-section arm, the targets reach a bending angle of $33.7^\circ$ and the trajectories $40.9^\circ$, compared with a 95th percentile of $35.6^\circ$ in the training data (Fig.~S9). The length change of the central axis is $-20.5$ to $-9.7$~mm for the single-section targets (training data, 5th to 95th percentile: $-24.4$ to $-7.0$~mm) and $-18.5$ to $-12.7$~mm per section for the multi-section targets (Tables~S8 and S9), so the evaluation covers a limited range of compression.
